\documentclass[a4paper]{scrartcl}
% \documentclass{article}

\usepackage[
fancytheorems, 
fancyproofs, 
noindent, 
%  spacingfix,  
]{adam}

\usepackage{tikz}
\usepackage{bbm}
\usepackage{mathtools}

\title{Topics in Combinatorics}
\author{Adam Kelly (\texttt{ak2316@cam.ac.uk})}
\date{\today}


\allowdisplaybreaks

\begin{document}

\maketitle

It is relatively hard to pin down exactly what Combinatorics is. Many different types of problems are studied, usually with some recurring characters: graphs, hypergraphs, sets of numbers, sets of points, subsets of $\{0, 1\}^n$, subsets of a group, and so on. These objects all have a `subset of a structured object' flavour, and it is this type of object that combinatorics usually deals with. In this article, we will be more concerned with various \emph{techniques} used in combinatorics.

This article constitutes my notes for the `Topics in Combinatorics' course, held in Lent 2021 at Cambridge and lectured by Prof. Tim Gowers. These notes are \emph{not a transcription of the lectures}, and differ significantly in quite a few areas. Still, all lectured material should be covered.


\tableofcontents

% \clearpage




\section{Averaging Arguments}

Let $X$ be an integrable real random variable: $\mathbb E|X|<\infty$. Then
\[
\mathbb P(X\geq\mathbb E X)>0.
\]
Indeed, otherwise $Y=\mathbb E X-X$ would be positive almost surely and integrable, with $\mathbb E Y=0$. Since $\{Y>0\}=\bigcup_{n\geq1}\{Y\geq1/n\}$ has probability one, some event in the union has positive probability. Thus $\mathbb E Y\geq\mathbb P(Y\geq1/n)/n>0$, a contradiction.
For a nonempty finite list of real numbers, this says that at least one entry is at least the arithmetic mean: if every entry were smaller, summing would contradict the definition of the mean. This is the averaging argument.

\end{document}